\chapter{2005年全国III卷（文）}
\section{单选题}
\begin{problem}
已知 \(\alpha\) 为第三象限角，则 \(\frac{\alpha}{2}\) 所在的象限是（\quad）

\choices
    {第一或第二象限}
    {第二或第三象限}
    {第一或第三象限}
    {第二或第四象限}
\end{problem}

\begin{answer}
D
\end{answer}

\begin{solution}
设 \(\alpha=2k\pi+\theta\)，其中 \(k\in\Z\)，\(\pi<\theta<\dfrac{3\pi}{2}\). 则
\(\dfrac{\alpha}{2}=k\pi+\dfrac{\theta}{2}\)，从而 \(\dfrac{\theta}{2}\in\left(\dfrac{\pi}{2},\dfrac{3\pi}{4}\right)\). 当 \(k\) 为偶数时，\(\dfrac{\alpha}{2}\) 在第二象限；当 \(k\) 为奇数时，\(\dfrac{\alpha}{2}\) 在第四象限，所以应选 D.
\end{solution}



\begin{problem}
已知过点 \(A(-2, m)\) 和 \(B(m, 4)\) 的直线与直线 \(2x + y - 1 = 0\) 平行，则 \(m\) 的值为（\quad）

\choices
    {\(0\)}
    {\(-8\)}
    {\(2\)}
    {\(10\)}
\end{problem}

\begin{answer}
B
\end{answer}

\begin{solution}
直线 \(2x+y-1=0\) 的斜率为 \(-2\). 当 \(m=-2\) 时，\(A(-2,-2)\) 与 \(B(-2,4)\) 所确定的直线为竖直直线，不可能与已知直线平行，因此 \(m\ne-2\). 于是
\(\dfrac{4-m}{m+2}=-2\)，解得 \(m=-8\)，所以应选 B.
\end{solution}



\begin{problem}
在 \((x - 1)(x + 1)^8\) 的展开式中 \(x^5\) 的系数是（\quad）

\choices
    {\(-14\)}
    {\(14\)}
    {\(-28\)}
    {\(28\)}
\end{problem}

\begin{answer}
B
\end{answer}

\begin{solution}
由二项式定理，\(x(x+1)^8\) 中 \(x^5\) 的系数为 \(\mathrm C_8^4\)，而 \(-(x+1)^8\) 中 \(x^5\) 的系数为 \(-\mathrm C_8^5\). 因此原展开式中 \(x^5\) 的系数为 \(\mathrm C_8^4-\mathrm C_8^5=70-56=14\)，所以应选 B.
\end{solution}



\begin{problem}
设三棱柱 \(ABC - A_{1}B_{1}C_{1}\) 的体积为 \(V\)，\(P\)，\(Q\) 分别是侧棱 \(AA_{1}\)，\(CC_{1}\) 上的点，且 \(PA = QC_{1}\)，则四棱锥 \(B - APQC\) 的体积为（\quad）

\choices
    {\(\frac{1}{6} V\)}
    {\(\frac{1}{4} V\)}
    {\(\frac{1}{3} V\)}
    {\(\frac{1}{2} V\)}
\end{problem}

\begin{answer}
C
\end{answer}

\begin{solution}
设 \(\overrightarrow{AA_1}=\overrightarrow{BB_1}=\overrightarrow{CC_1}=\bs u\)，并设 \(\overrightarrow{AP}=t\bs u\)，其中 \(0\leqslant t\leqslant1\). 设 \(\overrightarrow{CQ}=s\bs u\)，则由 \(PA=QC_1\) 得 \(t\lvert\bs u\rvert=(1-s)\lvert\bs u\rvert\)，所以 \(s=1-t\). 因而 \(AP\parallel CQ\)，且 \(AP+CQ=AA_1\).

四边形 \(APQC\) 位于由 \(\bs u\) 和 \(\overrightarrow{AC}\) 张成的平面内，以 \(AP\)、\(CQ\) 为平行边，其面积为
\(S_{APQC}=\dfrac12\lvert\bs u\times\overrightarrow{AC}\rvert\). 设点 \(B\) 到平面 \(APQC\) 的距离为 \(H\)，则
\(H=\dfrac{\left\lvert\overrightarrow{AB}\cdot\left(\bs u\times\overrightarrow{AC}\right)\right\rvert}{\lvert\bs u\times\overrightarrow{AC}\rvert}\). 因此四棱锥 \(B-APQC\) 的体积为
\(\dfrac13S_{APQC}H=\dfrac16\left\lvert\overrightarrow{AB}\cdot\left(\bs u\times\overrightarrow{AC}\right)\right\rvert\).

三棱柱的体积为 \(V=\dfrac12\left\lvert\overrightarrow{AB}\cdot\left(\overrightarrow{AC}\times\bs u\right)\right\rvert\)，而两个混合积的绝对值相等，所以四棱锥的体积为 \(\dfrac13V\)，应选 C.
\end{solution}



\begin{problem}
设 \(3^{x} = \frac{1}{7}\)，则（\quad）

\choices
    {\(-2 < x < -1\)}
    {\(-3 < x < -2\)}
    {\(-1 < x < 0\)}
    {\(0 < x < 1\)}
\end{problem}

\begin{answer}
A
\end{answer}

\begin{solution}
因为 \(3^{-2}=\dfrac19<\dfrac17<\dfrac13=3^{-1}\)，且函数 \(3^x\) 在 \(\R\) 上单调递增，所以 \(-2<x<-1\)，应选 A.
\end{solution}



\begin{problem}
若 \(a = \frac{\ln 2}{2}\)，\(b = \frac{\ln 3}{3}\)，\(c = \frac{\ln 5}{5}\)，则（\quad）

\choices
    {\(a < b < c\)}
    {\(c < b < a\)}
    {\(c < a < b\)}
    {\(b < a < c\)}
\end{problem}

\begin{answer}
C
\end{answer}

\begin{solution}
由 \(3\ln2<2\ln3\)（等价于 \(8<9\)），得 \(a<b\). 又由 \(2\ln5<5\ln2\)（等价于 \(25<32\)），得 \(c<a\). 因此 \(c<a<b\)，应选 C.
\end{solution}



\begin{problem}
设 \(0 \leqslant x \leqslant 2\pi\)，且 \(\sqrt{1 - \sin 2x} = \sin x - \cos x\)，则（\quad）

\choices
    {\(0 \leqslant x \leqslant \pi\)}
    {\(\frac{\pi}{4} \leqslant x \leqslant \frac{7\pi}{4}\)}
    {\(\frac{\pi}{4} \leqslant x \leqslant \frac{5\pi}{4}\)}
    {\(\frac{\pi}{2} \leqslant x \leqslant \frac{3\pi}{2}\)}
\end{problem}

\begin{answer}
C
\end{answer}

\begin{solution}
由 \(1-\sin2x=(\sin x-\cos x)^2\)，原方程等价于 \(\lvert\sin x-\cos x\rvert=\sin x-\cos x\)，所以必须有 \(\sin x-\cos x\geqslant0\). 于是 \(\sqrt2\sin\left(x-\dfrac\pi4\right)\geqslant0\). 在 \([0,2\pi]\) 上解得 \(\dfrac\pi4\leqslant x\leqslant\dfrac{5\pi}{4}\)，应选 C.
\end{solution}



\begin{problem}
\(\frac{2\sin 2\alpha}{1 + \cos 2\alpha} \cdot \frac{\cos^2\alpha}{\cos 2\alpha} =\)（\quad）

\choices
    {\(\tan \alpha\)}
    {\(\tan 2\alpha\)}
    {\(1\)}
    {\(\frac{1}{2}\)}
\end{problem}

\begin{answer}
B
\end{answer}

\begin{solution}
在原式有意义的范围内，利用 \(1+\cos2\alpha=2\cos^2\alpha\)，有 \(\dfrac{2\sin2\alpha}{1+\cos2\alpha}\cdot\dfrac{\cos^2\alpha}{\cos2\alpha}=\dfrac{\sin2\alpha}{\cos2\alpha}=\tan2\alpha\)，所以应选 B.
\end{solution}



\begin{problem}
已知双曲线 \(x^{2} - \frac{y^{2}}{2} = 1\) 的焦点为 \(F_{1}\)，\(F_{2}\)，点 \(M\) 在双曲线上且 \(\overrightarrow{MF_1} \cdot \overrightarrow{MF_2} = 0\)，则点 \(M\) 到 \(x\) 轴的距离为（\quad）

\choices
    {\(\frac{4}{3}\)}
    {\(\frac{5}{3}\)}
    {\(\frac{2\sqrt{3}}{3}\)}
    {\(\sqrt{3}\)}
\end{problem}

\begin{answer}
C
\end{answer}

\begin{solution}
由双曲线方程可知 \(a^2=1\)，\(b^2=2\)，所以 \(c^2=a^2+b^2=3\)，可取焦点为 \(F_1(-\sqrt3,0)\)、\(F_2(\sqrt3,0)\). 设 \(M(x,y)\)，则
\(\overrightarrow{MF_1}\cdot\overrightarrow{MF_2}=(-\sqrt3-x,-y)\cdot(\sqrt3-x,-y)=x^2+y^2-3=0\).
又因为 \(M\) 在双曲线上，所以 \(x^2-\dfrac{y^2}{2}=1\). 联立得 \(x^2=\dfrac53\)，\(y^2=\dfrac43\)，故点 \(M\) 到 \(x\) 轴的距离为 \(\lvert y\rvert=\dfrac{2\sqrt3}{3}\)，应选 C.
\end{solution}



\begin{problem}
设椭圆的两个焦点分别为 \(F_{1}\)，\(F_{2}\)，过 \(F_{2}\) 作椭圆长轴的垂线交椭圆于点 \(P\)，若 \(\triangle F_{1}PF_{2}\) 为等腰直角三角形，则椭圆的离心率是（\quad）

\choices
    {\(\frac{\sqrt{2}}{2}\)}
    {\(\frac{\sqrt{2} - 1}{2}\)}
    {\(2 - \sqrt{2}\)}
    {\(\sqrt{2} - 1\)}
\end{problem}

\begin{answer}
D
\end{answer}

\begin{solution}
设椭圆的长半轴、短半轴和半焦距分别为 \(a\)、\(b\)、\(c\)，则 \(a^2=b^2+c^2\). 取长轴为 \(x\) 轴，设 \(F_2=(c,0)\)，则过 \(F_2\) 的长轴垂线为 \(x=c\). 设其与椭圆的交点为 \(P\)，则
\(PF_2^2=b^2\left(1-\dfrac{c^2}{a^2}\right)=\dfrac{b^4}{a^2}\)，从而 \(PF_2=\dfrac{b^2}{a}\).

三角形 \(F_1PF_2\) 在 \(F_2\) 处为直角，且为等腰三角形，所以 \(PF_2=F_1F_2=2c\). 因此 \(\dfrac{b^2}{a}=2c\). 令离心率 \(e=\dfrac ca\)，代入 \(b^2=a^2-c^2\)，得 \(1-e^2=2e\). 因为 \(0<e<1\)，所以 \(e=\sqrt2-1\)，应选 D.
\end{solution}



\begin{problem}
不共面的四个定点到平面 \(\alpha\) 的距离都相等. 这样的平面 \(\alpha\) 共有（\quad）

\choices
    {\(3\text{ 个}\)}
    {\(4\text{ 个}\)}
    {\(6\text{ 个}\)}
    {\(7\text{ 个}\)}
\end{problem}

\begin{answer}
D
\end{answer}

\begin{solution}
把四个定点分成两个非空组，若平面 \(\alpha\) 到两组点的距离相等，则两组点分别位于 \(\alpha\) 两侧的两平行平面上，而 \(\alpha\) 是这两平行平面的中间平面.

按 \(1+3\) 分组时，任选一个点作为单独的一组，共有 \(4\) 种；按 \(2+2\) 分组时，把四点分成两对共有 \(\dfrac{\mathrm C_4^2}{2}=3\) 种. 因四点不共面，每一种分组都唯一确定相应的平行平面及其中间平面，所以共有 \(4+3=7\) 个，应选 D.
\end{solution}



\begin{problem}
计算机中常用十六进制是逢 \(16\) 进 \(1\) 的计数制，采用数字 \(0 \sim 9\) 和字母 \(\mathrm{A} \sim \mathrm{F}\) 共 \(16\) 个计数符号，这些符号与十进制的数的对应关系如下表：

\begin{center}
\small
\begin{tabular}{|c|c|c|c|c|c|c|c|c|}
\hline
十六进制 & \(0\) & \(1\) & \(2\) & \(3\) & \(4\) & \(5\) & \(6\) & \(7\) \\
\hline
十进制 & \(0\) & \(1\) & \(2\) & \(3\) & \(4\) & \(5\) & \(6\) & \(7\) \\
\hline
十六进制 & \(8\) & \(9\) & \(\mathrm{A}\) & \(\mathrm{B}\) & \(\mathrm{C}\) & \(\mathrm{D}\) & \(\mathrm{E}\) & \(\mathrm{F}\) \\
\hline
十进制 & \(8\) & \(9\) & \(10\) & \(11\) & \(12\) & \(13\) & \(14\) & \(15\) \\
\hline
\end{tabular}
\end{center}

例如，用十六进制表示：\(\mathrm{E} + \mathrm{D} = 1\mathrm{B}\)，则 \(\mathrm{A} \times \mathrm{B} =\)（\quad）

\choices
    {\(6\mathrm{E}\)}
    {\(72\)}
    {\(5\mathrm{F}\)}
    {\(\mathrm{B}0\)}
\end{problem}

\begin{answer}
A
\end{answer}

\begin{solution}
十六进制中 \(\mathrm A=10\)，\(\mathrm B=11\)，所以 \(\mathrm A\times\mathrm B=10\times11=110\). 又 \(110=6\times16+14\)，而十进制的 \(14\) 在十六进制中记为 \(\mathrm E\)，故 \(\mathrm A\times\mathrm B=6\mathrm E\)，应选 A.
\end{solution}


\section{填空题}
\begin{problem}
经问卷调查，某班学生对摄影分别执“喜欢”，“不喜欢”和“一般”三种态度，其中执“一般”态度的比“不喜欢”态度的多 \(12\text{ 人}\)，按分层抽样方法从全班选出部分学生座谈摄影. 如果选出的 \(5\text{ 位}\)“喜欢”摄影的同学，\(1\text{ 位}\)“不喜欢”摄影的同学和 \(3\text{ 位}\)执“一般”态度的同学，那么全班学生中“喜欢”摄影的比全班人数的一半还多\fillinblank{}\(\text{ 人}\).
\end{problem}

\begin{answer}
3
\end{answer}

\begin{solution}
设全班中“喜欢”、“不喜欢”和“一般”三种态度的学生人数分别为 \(5k\)、\(k\)、\(3k\). 由分层抽样的比例可知这三个数之比为 \(5:1:3\). 又由题意，\(3k-k=12\)，所以 \(k=6\). 因而全班人数为 \(5k+k+3k=54\)，喜欢摄影的学生人数为 \(5k=30\)，比全班人数的一半多 \(30-\dfrac{54}{2}=3\) 人.
\end{solution}



\begin{problem}
已知向量 \(\overrightarrow{OA} = (k, 12)\)，\(\overrightarrow{OB} = (4, 5)\)，\(\overrightarrow{OC} = (-k, 10)\)，且 \(A\)，\(B\)，\(C\) 三点共线，则 \(k =\)\fillinblank{}.
\end{problem}

\begin{answer}
\(-\frac{2}{3}\)
\end{answer}

\begin{solution}
由 \(\overrightarrow{AB}=(4-k,-7)\)，\(\overrightarrow{AC}=(-2k,-2)\)，且 \(A\)、\(B\)、\(C\) 三点共线，得
\((4-k)(-2)-(-7)(-2k)=0\). 化简得 \(-8-12k=0\)，所以 \(k=-\dfrac23\).
\end{solution}



\begin{problem}
曲线 \(y = 2x - x^3\) 在点 \((1,1)\) 处的切线方程为\fillinblank{}.
\end{problem}

\begin{answer}
\(x+y-2=0\)
\end{answer}

\begin{solution}
函数 \(y=2x-x^3\) 的导数为 \(y'=2-3x^2\)，所以在 \(x=1\) 处的切线斜率为 \(-1\). 过点 \((1,1)\) 的切线方程为 \(y-1=-(x-1)\)，即 \(x+y-2=0\).
\end{solution}



\begin{problem}
已知在 \(\triangle ABC\) 中，\(\angle ACB = 90^{\circ}\)，\(BC = 3\)，\(AC = 4\)，\(P\) 是 \(AB\) 上的点，则点 \(P\) 到 \(AC\)，\(BC\) 的距离乘积的最大值是\fillinblank{}.
\end{problem}

\begin{answer}
\(3\)
\end{answer}

\begin{solution}
建立平面直角坐标系，取 \(C=(0,0)\)，\(A=(4,0)\)，\(B=(0,3)\). 设 \(P=(x,y)\)，则 \(P\) 在线段 \(AB\) 上，且直线 \(AB\) 的方程为 \(3x+4y=12\)，其中 \(x\geqslant0\)、\(y\geqslant0\). 点 \(P\) 到 \(AC\)、\(BC\) 的距离分别为 \(y\)、\(x\)，故距离乘积为 \(xy\).

由 \(3x+4y=12\) 及基本不等式，\(12xy=(3x)(4y)\leqslant\left(\dfrac{3x+4y}{2}\right)^2=36\)，所以 \(xy\leqslant3\). 当且仅当 \(3x=4y\) 时取等号；联立得 \(x=2\)、\(y=\dfrac32\)，该点在线段 \(AB\) 上，故最大值为 \(3\).
\end{solution}



\section{解答题}
\begin{problem}
已知函数 \(f(x) = 2\sin^2 x + \sin 2x\)，\(x \in [0, 2\pi]\). 求使 \(f(x)\) 为正值的 \(x\) 的集合.
\end{problem}

\begin{solution}
由 \(\sin2x=2\sin x\cos x\)，得
\(f(x)=2\sin^2x+2\sin x\cos x=2\sin x(\sin x+\cos x)\).

在 \([0,2\pi]\) 上，\(\sin x=0\) 的零点为 \(0\)、\(\pi\)、\(2\pi\)，\(\sin x+\cos x=0\) 的零点为 \(\dfrac{3\pi}{4}\)、\(\dfrac{7\pi}{4}\). 逐段判断两因子的符号：在 \(\left(0,\dfrac{3\pi}{4}\right)\) 和 \(\left(\pi,\dfrac{7\pi}{4}\right)\) 内两因子同号，乘积为正；其余区间乘积为负或为零. 因此使 \(f(x)\) 为正值的 \(x\) 的集合为 \(\left(0,\dfrac{3\pi}{4}\right)\cup\left(\pi,\dfrac{7\pi}{4}\right)\).
\end{solution}


\begin{problem}
设甲，乙，丙三台机器是否需要照顾相互之间没有影响. 已知在某一小时内，甲，乙都需要照顾的概率为 \(0.05\)，甲，丙都需要照顾的概率为 \(0.1\)，乙，丙都需要照顾的概率为 \(0.125\).

\begin{enumerate}
\item 求甲，乙，丙每台机器在这个小时内需要照顾的概率分别是多少；
\item 计算这个小时内至少有一台需要照顾的概率.
\end{enumerate}
\end{problem}

\begin{solution}
\begin{enumerate}
\item 设甲、乙、丙三台机器需要照顾的概率分别为 \(p\)、\(q\)、\(r\). 由相互独立及题设，\(pq=0.05\)，\(pr=0.1\)，\(qr=0.125\). 因此
\(p^2=\dfrac{(pq)(pr)}{qr}=\dfrac{0.05\times0.1}{0.125}=0.04\)，从而 \(p=0.2\). 于是 \(q=\dfrac{0.05}{0.2}=0.25\)，\(r=\dfrac{0.1}{0.2}=0.5\).
\item 三台机器都不需要照顾的概率为 \((1-p)(1-q)(1-r)=0.8\times0.75\times0.5=0.3\). 所以至少有一台需要照顾的概率为 \(1-0.3=0.7\).
\end{enumerate}
\end{solution}


\begin{problem}
如图，在四棱锥 \(V - ABCD\) 中，底面 \(ABCD\) 是正方形，侧面 \(VAD\) 是正三角形，平面 \(VAD \perp\) 底面 \(ABCD\).

\begin{enumerate}
\item 证明：\(AB \perp\) 平面 \(VAD\)；
\item 求面 \(VAD\) 与面 \(VDB\) 所成的二面角的大小.
\end{enumerate}


\bitmapfigure[max width=0.62\linewidth]{img/2005/national_paper_3_liberal/q19_fig1.png}
\end{problem}

\begin{solution}
\begin{enumerate}
\item 设正方形边长为 \(s\)，建立空间直角坐标系，取 \(A=(0,0,0)\)、\(B=(s,0,0)\)、\(D=(0,s,0)\)、\(C=(s,s,0)\). 因为平面 \(VAD\) 垂直于底面，且 \(VAD\) 是边长为 \(s\) 的正三角形，可取 \(V=\left(0,\dfrac{s}{2},\dfrac{\sqrt3s}{2}\right)\). 于是平面 \(VAD\) 的方程为 \(x=0\)，而 \(AB\) 的方向向量为 \((s,0,0)\)，所以 \(AB\perp\) 平面 \(VAD\).
\item 取平面 \(VAD\) 的法向量 \(\bs{n}_1=(1,0,0)\). 平面 \(VDB\) 内的两条相交直线的方向向量可取 \(\overrightarrow{DB}=(s,-s,0)\)、\(\overrightarrow{DV}=\left(0,-\dfrac{s}{2},\dfrac{\sqrt3s}{2}\right)\)，故其法向量可取 \(\bs{n}_2=(\sqrt3,\sqrt3,1)\). 设两平面的锐二面角为 \(\theta\)，则
\(\cos\theta=\dfrac{\lvert\bs{n}_1\cdot\bs{n}_2\rvert}{\lvert\bs{n}_1\rvert\lvert\bs{n}_2\rvert}=\dfrac{\sqrt3}{\sqrt7}=\dfrac{\sqrt{21}}{7}\). 因而二面角的大小为 \(\arccos\dfrac{\sqrt{21}}{7}\).
\end{enumerate}
\end{solution}


\begin{problem}
在等差数列 \(\{a_{n}\}\) 中，公差 \(d \neq 0\)，\(a_{2}\) 是 \(a_{1}\) 与 \(a_{4}\) 的等差中项. 已知数列 \(a_{1}\)，\(a_{3}\)，\(a_{k_{1}}\)，\(a_{k_{2}}\)，\(\cdots\)，\(a_{k_{n}}\)，\(\cdots\) 成等比数列，求数列 \(\{k_{n}\}\) 的通项 \(k_{n}\).
\end{problem}

\begin{solution}
由等差数列通项公式，\(a_2=a_1+d\)，\(a_4=a_1+3d\). 题设“\(a_2\) 是 \(a_1\) 与 \(a_4\) 的等差中项”要求
\(2a_2=a_1+a_4\)，即 \(2(a_1+d)=a_1+(a_1+3d)\)，从而得到 \(d=0\)，这与题设 \(d\ne0\) 矛盾.

因此按题面字面条件，不存在满足题设的等差数列，不能推出数列 \(\{k_n\}\) 的通项.
\end{solution}


\begin{problem}
用长为 \(90 \mathrm{~cm}\)，宽为 \(48 \mathrm{~cm}\) 的长方形铁皮做一个无盖的容器，先在四角分别截去一个小正方形，然后把四边翻转 \(90^{\circ}\) 角，再焊接而成（如图），问该容器的高为多少时，容器的容积最大？最大容积是多少？


\FigureLayoutDeclare{img/2005/national_paper_3_liberal/q21_fig1.png}{12.0cm}{10bp 8bp 10bp 16bp}
\bitmapfigure[max width=0.70\linewidth]{img/2005/national_paper_3_liberal/q21_fig1.png}
\end{problem}

\begin{solution}
设截去的小正方形边长为 \(h\text{ cm}\)，则容器的高为 \(h\text{ cm}\)，且 \(0<h<24\). 容器底面的长、宽分别为 \(90-2h\text{ cm}\)、\(48-2h\text{ cm}\)，所以容积函数为
\(V(h)=h(90-2h)(48-2h)=4h^3-276h^2+4320h\).

求导得 \(V'(h)=12h^2-552h+4320=12(h-10)(h-36)\). 在 \(0<h<24\) 时，\(V'(h)>0\) 当 \(0<h<10\)，\(V'(h)<0\) 当 \(10<h<24\)，所以 \(h=10\) 时容积取得最大值. 此时
\(V(10)=10\times70\times28=19600\text{ cm}^3\). 因而容器的高为 \(10\text{ cm}\)，最大容积为 \(19600\text{ cm}^3\).
\end{solution}


\begin{problem}
设 \(A(x_{1}, y_{1})\)，\(B(x_{2}, y_{2})\) 两点在抛物线 \(y = 2x^{2}\) 上，\(l\) 是 \(AB\) 的垂直平分线.

\begin{enumerate}
\item 当且仅当 \(x_{1} + x_{2}\) 取何值时，直线 \(l\) 经过抛物线的焦点 \(F\)？证明你的结论；
\item 当 \(x_{1} = 1\)，\(x_{2} = -3\) 时，求直线 \(l\) 的方程.
\end{enumerate}
\end{problem}

\begin{solution}
\begin{enumerate}
\item 将抛物线写成 \(x^2=\dfrac12y\)，可知其焦点为 \(F\left(0,\dfrac14\right)\). 直线 \(l\) 是线段 \(AB\) 的垂直平分线，因此 \(F\) 在 \(l\) 上的充要条件是 \(FA=FB\). 由 \(A=(x_1,2x_1^2)\)、\(B=(x_2,2x_2^2)\)，有
\(FA^2-FB^2=4(x_1^2-x_2^2)(x_1^2+x_2^2)\).
其中 \(x_1^2+x_2^2>0\)，故 \(FA=FB\) 当且仅当 \(x_1^2=x_2^2\). 因为 \(A\)、\(B\) 为两点，\(x_1\ne x_2\)，所以这又等价于 \(x_1+x_2=0\). 因此，当且仅当 \(x_1+x_2=0\) 时，直线 \(l\) 经过焦点 \(F\).
\item 当 \(x_1=1\)、\(x_2=-3\) 时，\(A=(1,2)\)、\(B=(-3,18)\)，其中点为 \(M=(-1,10)\). 直线 \(AB\) 的斜率为 \(\dfrac{18-2}{-3-1}=-4\)，所以垂直平分线 \(l\) 的斜率为 \(\dfrac14\). 由点斜式得 \(y-10=\dfrac14(x+1)\)，即 \(l\) 的方程为 \(x-4y+41=0\).
\end{enumerate}
\end{solution}
